\section{The fixed coordinate model}
Write $\Omega=\F_{23}\cup\{\infty\}$, serializing $\infty$ as coordinate $23$. If $Q$ is the set of nonzero squares modulo $23$, let $\mathcal C\subset\F_2^{24}$ be the span of the 23 supports $(a+Q)\cup\{\infty\}$. A word is represented by the binary integer whose $i$th bit is its $i$th coordinate. The twelve words
\[
\begin{gathered}
\mathtt{85335e,
5606be,
2ca9ae,
11fe26,
0f55e2,
046f3e,}\\
\mathtt{02379f,
0149a6,
00a4d3,
0042de,
003e4a,
001f25}
\end{gathered}
\]
form a basis; the notation is hexadecimal. Row reduction of the defining supports gives this basis. Enumerating its $2^{12}$ sums gives weights $0,8,12,16,24$ with multiplicities $1,759,2576,759,1$. This is the projective Golay code \cite{chapman}. We use the Leech model \cite[Theorem 1.4]{brouwer}
\begin{equation}\label{P2:eq:model}
\La=\left\{x/\sqrt8:\begin{array}{l}x\in\Z^{24},\quad x_i\equiv m\pmod2\text{ for one }m\in\{0,1\},\\
\sum_i x_i\equiv4m\pmod8,\quad ((x_i-m)/2\bmod2)_i\in\mathcal C\end{array}\right\}.
\end{equation}
Thus $\La$ is even, unimodular and has minimum squared norm four. These properties and the lattice itself are inherited, not claimed as new.

The three disjoint octads are
\begin{align*}
B_0&=\{0,1,3,4,7,8,16,23\},\\
B_1&=\{10,12,15,18,19,20,21,22\},\\
B_2&=\{2,5,6,9,11,13,14,17\}.
\end{align*}
Put $w_i=2\one_{B_i}/\sqrt8$ and $P_i(x/\sqrt8)=\sum_{j\in B_i}x_j/4$. The matrices
\[
C=\begin{pmatrix}12&5\\1&12\end{pmatrix},\qquad
J=\begin{pmatrix}1&7\\3&22\end{pmatrix},\qquad
T=\begin{pmatrix}4&11\\16&4\end{pmatrix}
\]
act on $\Omega$ by fractional linear transformations over $\F_{23}$ and permute lattice coordinates. They preserve $\mathcal C$; this can be checked on its twelve basis words. In projective action $T^4=C$, and $C$ cycles $B_0,B_1,B_2$. All vector actions below are these coordinate permutations, not octonion-coordinate permutations.

\section{An octad-aligned coordinate isometry}
Use Wilson's octonion basis $(1,i_0,\ldots,i_6)$, with $i_0i_1=i_3$, its translates modulo seven, and anticommutativity. Thus $i_ti_{t+1}=i_{t+3}$ and each $i_t^2=-1$. Let $L\subset\mathbb O$ be the $E_8$ lattice generated by $\pm i_t\pm i_u$ and the half-sums having an odd number of minus signs; here $i_\infty=1$. Put
\[
s=(-1+i_0+\cdots+i_6)/2.
\]
Wilson's Leech lattice consists of triples $z=(z_0,z_1,z_2)$ with
\begin{equation}\label{P2:eq:wilson}
z_i\in L,\quad z_i+z_j\in L\bar s,\quad z_0+z_1+z_2\in Ls,
\qquad \|z\|_W^2=\tfrac12\sum_iN(z_i).
\end{equation}
The bar, side of multiplication and norm factor are part of the definition \cite[Sections 2--4]{wilson}.

\begin{theorem}[Specified model isometry]\label{P2:thm:isometry}
The following signed coordinate permutation identifies \eqref{P2:eq:model} with \eqref{P2:eq:wilson}, and aligns the three phase octads with the three octonions:
\[
\begin{array}{c|rrrrrrrr|rrrrrrrr}
j&\multicolumn{8}{c|}{\pi(8j),\ldots,\pi(8j+7)}&\multicolumn{8}{c}{\epsilon(8j),\ldots,\epsilon(8j+7)}\\\hline
0&0&1&3&4&8&23&7&16&+&-&-&-&-&-&-&-\\
1&10&22&15&18&12&21&19&20&-&-&-&-&+&+&-&+\\
2&2&13&5&11&6&14&17&9&-&-&-&-&+&+&-&+
\end{array}
\]
For $\lambda=x/\sqrt8$, its image has raw coordinates
\[
y_i=\epsilon_i x_{\pi(i)},\qquad
z_j=\frac12\sum_{r=0}^7y_{8j+r}e_r,
\quad(e_0,e_1,\ldots,e_7)=(1,i_0,\ldots,i_6).
\]
The inverse is $x_{\pi(i)}=\epsilon_i y_i$. Both maps are integral in their stated numerator coordinates and preserve squared norm.
\end{theorem}
\begin{proof}
The signed permutation is invertible and $\sum y_i^2/8=\sum x_i^2/8$, proving the metric statement. For lattice membership it suffices to test these 37 explicit generators of \eqref{P2:eq:model}:
\[
4(e_i-e_{23})\ (0\le i<23),\quad8e_{23},\quad
2c\ (c\text{ one of the twelve binary basis words}),\quad
(-3,1,\ldots,1).
\]
They are numerator vectors. To prove generation, the even part is generated by $4D_{24}$ and the twelve $2c$: binary addition differs from integer addition by four times an intersection vector, and those intersections have even weight because the code is self-orthogonal. The first 24 displayed vectors generate $4D_{24}$. Any two odd lattice vectors differ by an even one, so the final vector supplies the odd coset.

Here are exact integer tests for each image, permitting the 37 verifications without any geometric judgement. A block $v=2z$ represents a member of $L$ precisely when all entries have common parity $m$ and $\sum v_i\equiv2m\pmod4$. To test $z\in L\bar s$, multiply its raw block on the right by $(-1,1,\ldots,1)$; all product coordinates must be divisible by four and their quotients must pass the block test. To test $z\in Ls$, use $(-1,-1,\ldots,-1)$. These rules use $N(s)=2$ and retain all parentheses. Apply them to the three blocks, their three pair-sums and their total sum. The table above passes all 37 generator tests; the accompanying certificate lists every generator image and the independent checker performs the integer operations from the basis multiplication table.

This proves inclusion in Wilson's lattice. Wilson proves that \eqref{P2:eq:wilson} is unimodular at the displayed scale \cite[Sections 3--4]{wilson}. The source lattice is unimodular, so the isometric inclusion has index one. This upgrades generator checking to equality of the full infinite lattices; a minimal-shell census alone would not suffice.
\end{proof}
Wilson already supplies an octonionic-to-MOG comparison in Section 4. The assertion here is the displayed \emph{projective-octad-aligned} coordinate choice and its arithmetic readout, not the first existence of an isometry between Leech models.

\section{The exact Albert embedding and the retained cubic}
The three scalar traces become
\begin{equation}\label{P2:eq:traces}
P_j(\lambda)=\frac12\sum_{r=0}^7\epsilon_{8j+r}(z_j)_r
=\frac12\operatorname{Re}(z_j\overline{h_j}),\quad
h_j=\sum_{r=0}^7\epsilon_{8j+r}e_r.
\end{equation}
For the full, injective off-diagonal map put
\[
\iota(z)=\begin{pmatrix}0&z_2&\bar z_1\\\bar z_2&0&z_0\\z_1&\bar z_0&0\end{pmatrix}.
\]
Reading the marked entries is the inverse. The quadratic and cubic invariants are
\[
\operatorname{Tr}(\iota(z)^2)=4\|\lambda\|^2,\qquad
\mathcal N(\lambda)=N_J(\iota(z))=2\operatorname{Re}((z_0z_1)z_2).
\]
The association displayed in the last expression is fixed. In integer numerator coordinates,
\[
\mathcal N(\lambda)=\tfrac14\operatorname{Re}((y_0y_1)y_2),
\]
where $y_i$ now denotes the corresponding raw octonion block. Expansion has exactly 64 nonzero monomials; their complete signs and coordinate indices are serialized in the certificate.

\begin{proposition}[Two failed descents, with witnesses]\label{P2:prop:cubic}
The cubic $\mathcal N$ does not descend through $P$, and the specific phase isometry $C$ does not preserve $\mathcal N$.
\end{proposition}
\begin{proof}
The lattice numerator
\[
(-4,2,2,0,0,-2,-2,2,2,0,2,0,0,2,0,2,2,0,0,-2,-2,0,0,-4)
\]
has $P=0$ and $\mathcal N=-8$, whereas $\mathcal N(0)=0$. For
\[
b_0=\frac1{\sqrt8}(-3,1,-1,1,1,-1,-1,1,1,1,1,1,1,-1,1,1,1,1,-1,-1,-1,-1,1,1)
\]
one computes
\[
\mathcal N(b_0)=-5,\qquad\mathcal N(Cb_0)=1,\qquad\mathcal N(C^2b_0)=-1.
\]
Every assertion is direct lattice membership, a block-sum calculation, or multiplication in the stated octonion table.
\end{proof}
This does not prohibit a transported cubic: $g$ preserves the marked pair consisting of the lattice and the function $\mathcal N\circ g^{-1}$. It prohibits dropping that function's change of coordinates. The cyclic permutation of the three Wilson blocks is a different operator from the projective phase permutation $C$.

\section{A faithful combined state and its exact ES correction}
Let $t=P\lambda$, $n_i=N(z_i)$ and $\nu=2\operatorname{Re}((z_0z_1)z_2)$. Retain the full state
\[
\widehat X(\lambda)=\begin{pmatrix}t_0&z_2&\bar z_1\\\bar z_2&t_1&z_0\\z_1&\bar z_0&t_2\end{pmatrix}.
\]
Its image is specified by the lattice condition on $z$ and the three linear diagonal constraints \eqref{P2:eq:traces}; its inverse reads $z$ and applies Theorem~\ref{P2:thm:isometry}. The Albert formulas give
\[
N_J(\widehat X)=t_0t_1t_2-\sum_i t_in_i+\nu,
\qquad S_J(\widehat X)=\sum_{i<j}t_it_j-\sum_i n_i.
\]
Thus, for $F_p(t)=4t_0t_1t_2-p\sum_{i<j}t_it_j$,
\begin{equation}\label{P2:eq:fullcorrection}
\boxed{4N_J(\widehat X)-pS_J(\widehat X)
=F_p(t)+\sum_i(p-4t_i)n_i+4\nu.}
\end{equation}
Every mixed term is necessary. A norm-preserving transformation of the ambient Albert algebra is not automatically a map of this lattice graph or of the positive ES locus. Conversely, a lattice isometry need not preserve the fixed cubic, as Proposition~\ref{P2:prop:cubic} shows. Formula \eqref{P2:eq:fullcorrection} is the corrected interface, not a claim that these structures are unrelated.

\section{Verification and scope}
The independent standard-library checker verifies all 37 generator images and their inverses, and the compact-Zorn/three-block return over exact $\Q(i)$ coefficients. It tests the cubic identity on 2600 determining nodes: the 24 basis vectors, $e_i+e_j$ and $e_i-e_j$ for all pairs, and $e_i+e_j+e_k$ for all distinct triples. For a homogeneous cubic these separate every coefficient; testing only $e_i+e_j$ would not separate $x_i^2x_j$ from $x_ix_j^2$.

The note provides an explicit convention-compatible isometry and corrected polynomial interface. It asserts neither a new Leech lattice nor a solution of Erd\H{o}s--Straus. Content-level literature searches are recorded with the source package; historical priority for these particular coordinate formulas is not claimed. No theorem prover or independent reviewer was used.
